% ==============================================================================
% sci_s2_setup_sensors.tex — SCI/SICE tutorial Session 2:
% environment setup and sensor data acquisition
% SCI/SICE チュートリアル セッション2: 開発環境のセットアップとセンサデータの取得
% ==============================================================================

% ------------------------------------------------------------------
\begin{frame}{開発サイクルの地図と今の位置（セッション2）}
  \begin{columns}[c]
    \begin{column}{0.5\textwidth}
      \centering
      \resizebox{\columnwidth}{!}{\scisessionmap{2}}
    \end{column}
    \begin{column}{0.47\textwidth}
      \small
      セッション1で見た全体像のうち、まず実装の入口（環境構築とセンサ取得）から手を動かす
    \end{column}
  \end{columns}
\end{frame}

% ------------------------------------------------------------------
\begin{frame}{環境を整え、2 つの関数を書き、センサの値を見る}
  \begin{enumerate}
    \item ビルド環境は ESP-IDF と \code{sf CLI} が担う。導入も書き込みも sf CLI で行う
    \item 参加者が書くコードは \api{setup()} と \api{loop\_400Hz(dt)} の2関数だけ。センサやHWの詳細は \api{ws::} 名前空間が隠す
    \item センサデータの取得経路は3通りある。目的に応じて \api{ws::print}、\code{sf telemetry}、\code{sf log wifi} を使い分ける
  \end{enumerate}
\end{frame}

% ------------------------------------------------------------------
\begin{frame}[fragile]{OS による違いはここだけ}
  \begin{table}
    \centering\footnotesize
    \begin{tabular}{lll}
      \hline
      \textbf{OS} & \textbf{インストール} & \textbf{環境の有効化} \\
      \hline
      macOS / Linux  & \code{./install.sh}  & \code{source setup\_env.sh} \\
      Windows（CMD） & \code{install.bat}   & \code{setup\_env.bat} \\
      \hline
    \end{tabular}
  \end{table}

  \vspace{0.3em}

  {\footnotesize Windows は Git 同梱の CMD（コマンドプロンプト）でそのままファイル名を打つ。\code{./} も \code{source} も付けない}

  \vspace{0.5em}

  \begin{exampleblock}{書き込み時のシリアルポート名も違う}
    \footnotesize \code{sf flash vehicle -p /dev/ttyACM0}（Linux）／\code{-p /dev/cu.usbmodem*}（macOS）／\code{-p COM3}（Windows）
  \end{exampleblock}

  \vspace{0.3em}

  {\small 違いはこの2点だけ。\code{sf} コマンド自体（\code{sf build}、\code{sf lesson} 等）とワークフローは3\,OSで完全に同一}
\end{frame}

% ------------------------------------------------------------------
\begin{frame}{開発環境の導入: 機体とペアリング}
  \begin{block}{書き込みは \texttt{sf flash vehicle} / \texttt{sf flash controller} で行う}
    本講座の操作はすべて \texttt{sf} CLI で行う
  \end{block}

  \vspace{0.5em}

  \begin{exampleblock}{コントローラのペアリング（初回のみ）}
    \footnotesize
    \begin{enumerate}\setlength{\itemsep}{2pt}
      \item コントローラの LCD パネルボタンを押しながら電源を入れる
      \item StampFly 本体のボタンを \warn{3 秒以上}押し続け、ビープが鳴ったら離す（5 秒以上押し続けるとシステムリセット）
      \item 双方がビープしたらペアリング完了。以降は電源を入れるだけで自動再接続する
    \end{enumerate}
  \end{exampleblock}
\end{frame}

% ------------------------------------------------------------------
\begin{frame}[fragile]{実習 1 (1/2): 環境確認とビルド}
  \vspace{0.15em}
  {\small 見るだけ ｜ \textcolor{sfblue}{\textbf{一緒に打つ}} ｜ 帰宅後に再現}

  \begin{block}{手順}
    \begin{enumerate}\footnotesize\setlength{\itemsep}{0pt}
      \item \code{setup\_env.bat}（macOS / Linux は \code{source setup\_env.sh}）
      \item \code{sf doctor}（ESP-IDF・Python・USBドライバ・sf CLI 自体を診断）
      \item \code{sf build vehicle}（機体ファームをビルド。初回は数分かかる）
      \item \code{sf flash vehicle -m}（USB で書き込み、続けてモニタを開く。終了は \code{Ctrl+]}）
    \end{enumerate}
  \end{block}

  \begin{exampleblock}{確認すること}
    \footnotesize 書き込み直後に起動音が鳴り、LED が白（初期化中・静置）$\to$ 緑常灯（準備完了）になる。モニタに起動ログが流れる
  \end{exampleblock}

  {\footnotesize \code{sf doctor} が全て OK にならなくても心配は不要。このあとは「見るだけ」で参加し、復習パス（本セッション末尾）に沿って会場外で環境を整えてほしい}
\end{frame}

% ------------------------------------------------------------------
\begin{frame}[fragile]{実習 1 (2/2): 機体の初期設定を一度に済ませる}
  \vspace{0.2em}
  \begin{columns}[T]
    \begin{column}{0.50\textwidth}
      \begin{block}{USB をつないだまま、モニタ（CLI）で順に打つ}
        \footnotesize 機体ごとに 1 回。N は講師が指定（1 / 6 / 11）
        \begin{sflisting}[basicstyle=\ttfamily\scriptsize]{sf monitor の CLI}
param reset
param set wifi.mode 1
param set wifi.channel N
param save
reboot
        \end{sflisting}
      \end{block}
    \end{column}
    \begin{column}{0.46\textwidth}
      \begin{exampleblock}{続けてコントローラとペアリング}
        \footnotesize コントローラの LCD パネルボタンを押しながら電源を入れ、機体のボタンを \warn{3 秒以上}押し続ける。双方がビープしたら離す（\warn{5 秒以上}押し続けるとシステムリセット）。以降は電源投入だけで自動再接続
      \end{exampleblock}
      {\footnotesize \code{wifi.mode 1}: 機体が自分の WiFi を出す（後半で使う）。\code{wifi.channel}: 混信を避ける機体ごとのチャンネル。コントローラはペアリング時に自動で合わせる}
    \end{column}
  \end{columns}
\end{frame}

% ------------------------------------------------------------------
\begin{frame}{ツールチェーン}
  \begin{block}{ESP-IDF の上に sf CLI を重ねる}
    \small \code{idf.py} を直接叩く代わりに、ビルド・書き込み・診断・ログ取得を統一したコマンド体系で提供する
  \end{block}

  \vspace{0.5em}

  \begin{table}
    \centering\footnotesize
    \begin{tabular}{ll}
      \hline
      \textbf{コマンド} & \textbf{役割} \\
      \hline
      \code{sf build [target]} & ファームウェアビルド \\
      \code{sf flash [target] -m} & 書き込み（\code{-m} でモニタ付き） \\
      \code{sf monitor}          & シリアルモニタ \\
      \code{sf doctor}           & 環境診断 \\
      \hline
    \end{tabular}
  \end{table}

  \vspace{0.5em}

  {\small \code{target} には \code{vehicle}（機体）/ \code{controller}（送信機）を指定する。実習 1 で打った 4 コマンドがこの表の全てである}
\end{frame}

% ------------------------------------------------------------------
\begin{frame}[fragile]{参加者が書くコードと作業の流れ}
  {\small 参加者が編集するファイルは \code{user\_code.cpp} の1つだけ。本資料ではこれを「実習コード」と呼ぶ}
  \vspace{0.2em}
  \begin{columns}[T]
    \begin{column}{0.56\textwidth}
      \begin{block}{参加者が書く関数は2つだけ}
        \footnotesize
        \begin{itemize}
          \item \api{setup()} --- 起動時に1回
          \item \api{loop\_400Hz(dt)} --- 400\,Hz毎
        \end{itemize}
      \end{block}
      \vspace{0.3em}
      \begin{exampleblock}{切替・編集・ビルド・書き込み}
        \footnotesize
        \code{sf lesson switch sci2026:N}\\
        \code{sf lesson edit}\\
        \code{sf lesson build}\\
        \code{sf lesson flash}\\
        \code{sf lesson monitor}
      \end{exampleblock}
    \end{column}
    \begin{column}{0.40\textwidth}
      \centering
      \adjustbox{max width=\columnwidth, max height=0.72\textheight}{\input{../../_shared/tikz/sci_build_flash_flow}}
    \end{column}
  \end{columns}
\end{frame}

% ------------------------------------------------------------------
\begin{frame}{機体の軸定義}
  \vspace{-0.5em}
  \begin{center}
    % Width and height both capped (BL10): the figure is a true 3D projection now,
    % so its aspect ratio is no longer tied to the slide's
    % 幅と高さの両方を上限で制限（BL10）。図は正しい 3D 投影になり縦横比が変わりうるため
    \adjustbox{max width=0.62\textwidth, max height=0.52\textheight}{\input{../../_shared/tikz/imu_axes}}
  \end{center}
  \vspace{-0.5em}
  \begin{block}{}
    \centering\small
    BMI270 --- 6軸（3軸ジャイロ + 3軸加速度）400\,Hz \quad|\quad \mbox{Body FRD（前方-右-下）座標系}
  \end{block}
  {\footnotesize 機体の Body 座標系（前方-右-下）を先に決め、IMU の軸はそれに一致するように扱っている}
\end{frame}

% ------------------------------------------------------------------
\begin{frame}{IMU API と単位}
  \begin{table}
    \centering\footnotesize
    \begin{tabular}{lll}
      \hline
      \textbf{関数} & \textbf{説明} & \textbf{単位} \\
      \hline
      \api{gyro\_x/y/z()}  & 角速度（Roll/Pitch/Yaw） & rad/s \\
      \api{accel\_x/y/z()} & 加速度（X/Y/Z） & m/s\textsuperscript{2} \\
      \hline
    \end{tabular}
  \end{table}

  \vspace{0.5em}

  \begin{exampleblock}{静止時の目安}
    \small ジャイロ $\approx 0$、加速度 Z $\approx -9.81$\,m/s\textsuperscript{2}（重力に対する反力）
  \end{exampleblock}
\end{frame}

% ------------------------------------------------------------------
\begin{frame}[fragile]{実習 2: IMU の値を見る}
  \vspace{0.15em}
  {\small 見るだけ ｜ \textcolor{sfblue}{\textbf{一緒に打つ}} ｜ 帰宅後に再現}

  \begin{block}{手順}
    \begin{enumerate}\footnotesize\setlength{\itemsep}{1pt}
      \item \code{sf lesson switch sci2026:2}
      \item \code{sf lesson edit} で \code{user\_code.cpp} を開き、次ページのコードを書いて保存
      \item \code{sf lesson build}
      \item \code{sf lesson flash}
      \item \code{sf lesson monitor}（\code{sf lesson flash} 直後はそのまま開く）
    \end{enumerate}
  \end{block}

  \begin{exampleblock}{確認: Teleplot の波形}
    \small VSCode 拡張 \code{alexnesnes.teleplot} 上で \api{gyro\_x}/\api{gyro\_y} の波形が動くこと
  \end{exampleblock}

  {\footnotesize 帰宅後に再現: 同じ手順（switch $\to$ edit $\to$ build $\to$ flash $\to$ monitor）をそのまま実行すればよい}
\end{frame}

% ------------------------------------------------------------------
\begin{frame}[fragile]{データ取得経路① Teleplot}
  \begin{block}{\texttt{ws::print} をそのままグラフにする}
    \small \code{>変数名:値} の形式で出力すると、VSCode拡張機能 Teleplot がリアルタイムにグラフ化する
  \end{block}

  \begin{sflisting}[basicstyle=\ttfamily\scriptsize]{user\_code.cpp}
static uint32_t tick = 0;
void loop_400Hz(float dt) {
    if (tick++ % 4 == 0) {   // 100Hz decimation
        ws::print(">gyro_x:%.3f", ws::gyro_x());
        ws::print(">gyro_y:%.3f", ws::gyro_y());
    }
}
  \end{sflisting}

  \vspace{-4.5mm}

  {\footnotesize VSCode 拡張 \code{alexnesnes.teleplot} を導入後、\code{sf lesson monitor} で接続（\code{sf lesson flash} 直後はそのまま開く）}
\end{frame}

% ------------------------------------------------------------------
\begin{frame}[fragile]{実習 2: 傾けて確認}
  \vspace{0.15em}
  {\small 見るだけ ｜ \textcolor{sfblue}{\textbf{一緒に打つ}} ｜ 帰宅後に再現}

  \vspace{0.3em}

  \begin{block}{何を見るか}
    \begin{itemize}\small
      \item StampFly を手に持ち、前後左右に傾ける
      \item \api{gyro\_x}/\api{gyro\_y} が傾ける速さに応じて振れる
      \item 静止させると \api{accel\_z} が $-9.81$ 付近に戻る
    \end{itemize}
  \end{block}

  {\footnotesize 前ページまでの \code{sf lesson build \&\& sf lesson flash} 完了後、Teleplot 画面を見ながら機体を傾ける}
\end{frame}

% ------------------------------------------------------------------
\begin{frame}[fragile]{WiFi でつなぐ: 有線と無線の使い分け}
  \begin{exampleblock}{飛行中は USB を挿せない}
    \footnotesize USB ケーブルは機体と PC を物理的につなぐため、実際に飛ばしながらは使えない。
    その場でのゲイン変更や飛行中のデータ取得には WiFi 接続が要る
  \end{exampleblock}

  \vspace{0.3em}

  \begin{table}
    \centering\footnotesize
    \begin{tabular}{lcc}
      \hline
       & \textbf{USB（有線）} & \textbf{WiFi（無線）} \\
      \hline
      ビルド・書き込み                   & \checkmark & --- \\
      起動ログ確認（\code{sf monitor}）   & \checkmark & --- \\
      50\,Hz テレメトリのライブ表示        & --- & \checkmark \\
      飛行中の 400\,Hz 全データ記録        & --- & \checkmark \\
      パラメータの即時変更（飛行中も反映） & --- & \checkmark \\
      システム同定・ゲイン設計             & --- & \checkmark \\
      \hline
    \end{tabular}
  \end{table}

  {\footnotesize 右列は \code{sf telemetry}（UDP:5005）／\code{sf log wifi}（UDP:8890）／
  \code{param set}（WiFi 経由の CLI，TCP:23）のいずれかを使う。次ページで接続方法を説明する}
\end{frame}

% ------------------------------------------------------------------
\begin{frame}[fragile]{接続手順: 機体を WiFi 網に入れる}
  \begin{block}{機体側: 実習 1 (2/2) の初期設定で済んでいる}
    \footnotesize \code{wifi.mode 1}（機体が自分の WiFi を出す設定）は初期設定で保存済み。未設定の機体だけ、USB 接続のまま \code{sf monitor} の CLI で \code{param set wifi.mode 1} $\to$ \code{param save} $\to$ \code{reboot}
  \end{block}

  \vspace{0.15em}

  \begin{exampleblock}{PC 側: 機体の WiFi に接続する}
    \footnotesize 機体が \texttt{StampFly-XXYY}（MAC 末尾から自動生成の SSID）を出す。
    PC の WiFi 設定で選び，既定パスワード \texttt{stampfly} を入力。
    IP は既定 \texttt{192.168.10.1}（DJI Tello 互換サブネット）
  \end{exampleblock}

  \vspace{0.15em}

  {\scriptsize 接続の確認: \code{sf telemetry}（受信できれば OK，IP 指定不要）。
  会場の WiFi ではなく機体自身が出す AP に繋ぐ（機体ごとに SSID が異なる）}
\end{frame}

% ------------------------------------------------------------------
\begin{frame}[fragile]{データ取得経路② sf telemetry\later}
  \begin{block}{コード変更なしでライブ表示}
    \small 機体は状態推定値を常時 50\,Hz で送信している。受信して表示するだけなら実習コードの変更は不要
  \end{block}

  \begin{lstlisting}[style=sfbash, numbers=none]
sf telemetry          # ターミナル表示（既定）
sf telemetry --web    # ブラウザ表示（UDP:5005 -> SSE）
  \end{lstlisting}

  \vspace{0.3em}

  {\small 講義や複数人での画面共有には \code{--web} が見やすい。UDP:5005 を SSE（サーバーからブラウザへの一方向配信）でブラウザへ届けている}
\end{frame}

% ------------------------------------------------------------------
\begin{frame}[fragile]{データ取得経路③ ログ取得と可視化\later}
  \begin{block}{WiFi 経由で 400\,Hz Data Stream を記録する}
    \small \code{sf log wifi} は機体と同じ WiFi（SoftAP モードなら機体が出す AP）で 400\,Hz の Data Stream を UDP 受信し、フライトログ一式（\code{.sflog.zip}：信号ごとの CSV をまとめた zip 1個，埋め値なし）として保存する
  \end{block}

  \begin{lstlisting}[style=sfbash, numbers=none]
sf log wifi -d 30    # 30秒キャプチャ -> logs/flight_<日時>.sflog.zip
sf log viz           # 既定で logs/ 内の最新の一式を描画
  \end{lstlisting}

  \vspace{0.3em}

  \begin{exampleblock}{使い分けの目安}
    \footnotesize Teleplot = その場で波形を見る／\code{sf telemetry} = コード変更なしで確認／\code{sf log wifi} = 後で解析するためのデータ保存
  \end{exampleblock}
\end{frame}

% ------------------------------------------------------------------
\begin{frame}{他のセンサ\later}
  \small IMU 以外のセンサも同じ考え方で読める。時間の都合で本セッションでは扱わないが、単独ビルド可能な例題として用意してある

  \vspace{0.5em}

  \begin{table}
    \centering\footnotesize
    \begin{tabular}{ll}
      \hline
      \textbf{例題} & \textbf{センサ} \\
      \hline
      \code{firmware/vehicle/examples/05\_read\_tof} & VL53L3CX（ToF距離センサ） \\
      \code{firmware/vehicle/examples/06\_read\_baro} & BMP280（気圧センサ） \\
      \hline
    \end{tabular}
  \end{table}

  \vspace{0.5em}

  {\footnotesize 実習 1 と同じ sf CLI で単独ビルド・書き込みできる: \code{sf build vehicle/examples/05\_read\_tof} $\to$ \code{sf flash vehicle/examples/05\_read\_tof -m}}
\end{frame}

% ------------------------------------------------------------------
\begin{frame}{理論とコードの対応表\later}
  {\footnotesize 目的: 理論の各項目が実装のどこにあるか（関数・Data Stream 列・トピック）を後で引ける索引にする}

  \vspace{0.3em}

  \begin{table}
    \centering\footnotesize
    \begin{tabular}{lccc}
      \hline
      \textbf{理論} & \textbf{コード} & \textbf{Data Stream 列} & \textbf{トピック} \\
      \hline
      角速度 $\omega$ [rad/s]、Body FRD & \api{ws::gyro\_x()} & \code{gyro\_x} & \code{sensor\_imu} \\
      並進加速度 $a$ [m/s\textsuperscript{2}] & \api{ws::accel\_x()} & \code{accel\_x} & \code{sensor\_imu} \\
      \hline
    \end{tabular}
  \end{table}

  \vspace{0.5em}

  {\small 参加者が呼ぶ \api{ws::} 関数は、内部で \code{sensor\_imu} トピックを購読しているだけである
  （アクセス階層 L0（実習用 API）$\to$ L1（Topic API）の対応）}
\end{frame}

% ------------------------------------------------------------------
\begin{frame}{復習パス\later}
  \begin{table}
    \centering\footnotesize
    \setlength{\tabcolsep}{4pt}
    \begin{tabular}{>{\raggedright\arraybackslash}p{50mm}>{\raggedright\arraybackslash}p{28mm}>{\raggedright\arraybackslash}p{55mm}}
      \hline
      \textbf{コマンド} & \textbf{実習} & \textbf{文書} \\
      \hline
      \code{sf doctor} / \code{sf build vehicle} / \code{sf flash vehicle -m} & 実習1（環境確認とビルド） & \code{docs/\allowbreak guides/\allowbreak troubleshooting.md} \\
      \code{sf lesson switch sci2026:2} & 実習2（IMU） & \code{firmware/\allowbreak vehicle/\allowbreak docs/\allowbreak architecture.md \S2} \\
      \code{sf log wifi} / \code{sf log viz} & --- & \code{docs/\allowbreak guides/\allowbreak flight-log-viz.md} \\
      \hline
    \end{tabular}
  \end{table}
\end{frame}

% ------------------------------------------------------------------
\begin{frame}{チェックポイント}
  \begin{alertblock}{確認事項}
    \begin{itemize}
      \item[$\square$] \code{sf doctor} が通る
      \item[$\square$] \api{setup()}/\api{loop\_400Hz(dt)} の2関数構造を説明できる
      \item[$\square$] ジャイロ・加速度をどれか1つの経路で確認できた
    \end{itemize}
  \end{alertblock}

  \vspace{1em}

  \begin{block}{次のセッション}
    セッション3: モータ制御とコントローラ入力の実装 --- センサの次は、機体を動かす側を作る
  \end{block}
\end{frame}
